\chapter{积极心态}
相信大家既然历经重重考验来到了北京大学工学院，心理素质都是极强的，每个人都有不同的办法调节心情心态。在军训及新生训练营期间，同学们也会接触到一系列心理讲座，让大家对北大的学习和生活建立良好的心态。在这里，我们针对性地列出一些工学院学生常见的烦恼，并提供一些建议，希望能帮助同学们做好基本心理建设，以积极向上的姿态投入到学习和生活之中。当然，建设强大的内心本就是一个见仁见智的过程，我们希望大家能够吸取其中对自己有用的建议，大胆摒弃行之无效的方法，走自己的道路，过自己的美好人生！

\section{新生面临的核心问题}
这是一个很宽泛同时又很重要的话题，笔者不敢妄下结论，抛砖引玉，供管中窥豹。

\vspace{10pt}

大学生活与高中生活的不同，要从日常生活和工作两方面说起。

\vspace{10pt}

在日常生活方面，大学可能是很多同学第一次长时间离开家，随另一个群体生活。与高中生活不同的是，我们得不到高中那样不请自来的关注和指导，因此绝大多数事情需要我们自己去问、自己去摸索，而不能依赖别人主动地悉数灌输给我们。这时候，\textbf{独立性}和\textbf{主动性}就显得格外重要。我们需要学会自己探索，主动求教，为自己的生活负责，这样才可以克服生活中遇到的种种难题。

\vspace{10pt}

此外，我们需要面临更真切、更近距离的来自人际关系方面的压力。室友，同学，辅导员，老师……处理人际关系的责任不容推脱地落在了我们身上。而由于不同的文化环境、性格、目标等因素，孤独、矛盾等是不可避免的。笔者想说：人际关系最终的处理水平，仍然取决于我们的主观选择。如果他人没有恶意，宽容和体谅总不会出错。

\vspace{10pt}

可能有同学已经注意到笔者分类模式是日常生活和“工作”而非“学习”，这是因为人各有志，一些同学不将学术作为未来人生的总目标，这是完全可以理解的。因此我们广义地讲“工作”。

\vspace{10pt}

在工作方面，同学们会感觉到前所未有的自由，因为没有人会实时紧盯你的状态，工作和娱乐的权利全在你的手上。如果说高中是轨道，那么大学便是旷野；高中是一幅不容一点瑕疵的工笔画，那么大学便是一张一尘不染的白纸。白纸意味着没有上限，当然也没有下限；这句的前半句是激励，后半句是警告。没有人会阻止我们过放浪形骸的生活，当然也没有人会为我们设定天花板。北大便是这样，资源就摆在那里，不会主动喂到你的嘴边；但是只要你愿意争取，那学校便毫不吝啬地把资源向你倾斜。也许，一些同学已经对“自律”一词有了过犹不及的嫌恶，但是这里笔者仍然将\textbf{“自律”}作为一条极其重要的建议。当然，我们不鼓励打鸡血式地内卷，而是希望同学们想清楚自己想要什么、想成为什么样的人，并实时审视自己是离目标越来越近还是渐行渐远。如果适量的惬意休息是海船暂时归港，那么不易察觉的长期放纵便是摧毁海船的暗流与漩涡。

\vspace{10pt}

总之，建议有三：一是为自己负责，二是独立与主动，三是宽容与体谅。笔者真诚的期待同学们可以迎来四年无怨无悔的大学生活。

\section{如何利用好资源寻求心理帮助}
学校为我们开设了完备的心理服务。心理咨询中心会开设各类主题活动，可以关注公众号“北大学生心理健康教育与咨询中心”以参加相应活动。如果遇到难以解决的困难，可以拨打心理咨询中心电话：前台62760852 ；24小时心理援助热线62760521。

\vspace{10pt}

我们应该始终记得：心理障碍或创伤，和肉体上的创伤是一样的，是不可抗因素，不显示主观意志的强弱。当我们觉得难以支撑、心力交瘁时，要知道，这不是因为我们脆弱或敏感。因此，请积极地向外界寻求帮助，这并不羞耻。相信心理障碍可以被克服，我们终究会迎来拨云见日之时。

\section{常见问题荟萃}

说明：本部分内容仍待拓展，欢迎大家在我们的问卷中留言或投稿，也可以向我们私发消息。我们会选取具有普遍性的问题进行研究和约稿后写入手册（问卷二维码见手册“后记”）。

\subsection{学术类}
\textbf{（1）没有竞赛基础怎么适应工学院的学习生活？}

相对其他学院，工学院的竞赛基础没有那么夸张，请大家放心“祛魅”。且工学院相对更重视计算而非技巧，因此不必因为无竞赛背景而感到压力。此外，大学课程的重点与高中竞赛并非完全重合，有竞赛基础的同学也不能过于掉以轻心。

\vspace{10pt}

\textbf{（2）选了数院/物院/信科的课程成绩很差怎么办？}

首先，我们要承认其他人在天赋和基础上的优势，再看自己是否能够在意志品质上找补。如果发现成绩不尽如人意，那么建议综合权衡得分和难度，再进行选择。诚然，数院物院的课程相对难度更大，抽象程度更高，但不见得就不适合工学院同学进行学习\footnote{毕竟力学系学生拿过数学系数分第一名的传说不是盖的}，具体如何选择需要同学们自己衡量。

\vspace{10pt}

另外，在一门课程中取得不理想的绩点，并不意味着“天塌了”。一门专业课的成绩不佳也不能说明你的学习能力不如他人。或许是由于对课程难度的预估不准，或许是由于某次的考试失利。这一次失利放在整个大学上百学分的课程的学习之中也不过沧海一粟。同学们不必对自己过分苛刻，调整心态，再接再厉即可。

\vspace{10pt}

\textbf{（3）身边人都开始卷绩点/科研，我怎么办？}

生活不是竞技场。看别人干什么就跟风干什么，最终一定是内卷到疲惫且痛苦的。大学并非单行道，“3.92”\footnote{指绩点，是一个树洞里的梗}也不是所有问题的通解。因此，我们最好在一开始就做好大概预期，未来是做学生工作，还是搞科研，是留在国内还是去国外，然后再决定自己的生活节奏。不要人云亦云，做无意义的焦虑。

\vspace{10pt}

当然，也不建议因为没想好上述事情而焦虑，关于前途的考量，很多人是大三大四才定下来的，此所谓“车到山前必有路，船到桥头自然直”的智慧。笔者认为，大一最主要的任务无非就两件事：适应大学节奏、打好专业基础。真正做好这两件事就很好了。

\vspace{10pt}

至于科研，大一的同学更不必过早的焦虑，学好专业基础课后，再慢慢了解各个课题组，发掘自己的感兴趣的方向即可，“为了科研而科研”是行不通的。（夹带私货：有强烈的意愿的同学也可以关注院学生会举办的本研动力学讲座~）

\textbf{（4）面对大一末的分流，我应该做什么准备？}

在大一期间，同学们可以通过多种方式了解工学院的各个方向。大一上的现代工学通论课程，大一下的“工学嘉年华”活动、“工学文化节”，以及在分流前的专业介绍讲座，都为同学们提供了多样的了解工学院不同专业的方式。分流截止至大一结束的暑假初期，在大一下期末会进行一次预选以及召开一次针对预选结果的年级大会。大家也可以多与不同方向的学长学姐沟通交流，充分了解后进行选择。

\subsection{生活类}
\textbf{（1）除了做题，感觉自己其他啥也不会怎么办？}

\vspace{10pt}

工学院某位著名的“人生导师”曾经说过：“在北大，千万不要把自己太当回事，也千万不要把自己不当回事；千万不要把别人太当回事，也千万不要把别人不当回事。”无论是感觉自己“啥也不会”，还是感觉别人“啥都会”，事实上都是太把别人当回事，却把自己不当回事了。事实上，北大将教会你的第一件事情就是“祛魅”。让身边许许多多人不再被“神化”，许许多多事情不再“遥不可及”，大家都有自己的优点和缺点，很多也都是经历了高中的“小镇做题”阶段才进入北大的\footnote{必须承认的是，会做题本身也是很大的优点}。大家都是一张白纸，等待着自己花四年甚至更多时间，书写在北大的精彩人生。

\vspace{10pt}

北大会让你认识到的第二件事情，就是有梦想就会有实现的机会。第4章已经说了，只要想得到的活动一般都能找得到。在北大，资源不会主动找上门。但如果你积极地去找，一般都能找到很好的资源进行学习，同时试错的成本极低\footnote{花钱少，不用找关系，地理位置近，教学内容专精……}。因此，即便你现在觉得自己“啥也不会”，等过了一个学期、一年、四年，你会发现自己冥冥之中有了巨大的提升。

\vspace{10pt}

第三，数院有人说过：“来北大我可能没有成为最好的那个，但至少我知道了最好的是什么样子。”即便发现自己有些事情努力了也做不成，但是你已经去尝试过，也见过最好的事物是怎么样的，这本身就是一个眼界和能力的巨大提升，对未来帮助很大。同时，在未来你仍然有机会再次挑战它。记住，你才18岁。

\vspace{10pt}

\textbf{（2）学工、活动占据过多时间，导致成绩下降怎么办？}

\vspace{10pt}

工学院的一位著名学长曾经说过：“大学就是在不同赛道上成就自我。”大学赛道不再单一化，成绩不再是成功的唯一衡量标准。在成绩基本盘hold住的情况下\footnote{指至少不能挂科，成绩还说得过去}，有的同学选择深入钻研课程，力争更高的成绩，为将来保研出国和深入研究打下基础；有的同学选择将大部分时间投身科研，追求知识的深度，早早开始挑战真正有希望为人类世界带来GDP的业界技术难题；有的同学选择学习辅双或选择其他院系的课程，拓展自己的学术兴趣、志趣，在知识的广度上实现自我价值；有的同学选择投入学工，发挥自己的社交属性，和老师、同学们深入交流，并在工作中锤炼实用技能，为将来的工作能力奠基；有的同学不断寻找实习，为将来进入业界打下坚实的经验基础，同时凭着实习的工资实现小小的“财务自由”；有的同学选择多多参与文体活动，在活动中绽放自我，结交朋友，成为舞台上最闪亮的星星；有的人拓展自己的爱好，坚持养成一个个好习惯，不在外表上宣扬自己的强大，内核却永远在悄悄变强，期待着惊艳大家的那一天……以上种种，都是“赛道思维”很好的体现。只要是实现了自己的价值就有意义，不能说哪条赛道就一定比另一条好，正所谓“不管黑猫白猫，能抓到老鼠的就是好猫。”

\vspace{10pt}

因此，如果出现了上述类似的问题，笔者的建议是：你需要衡量自己心目中各条赛道的优先级，坚守课业成绩基本盘不动摇的同时，摒弃“成绩为王”的高中生思维定势，选择最能实现自我的道路坚定地发展。

\vspace{10pt}

\textbf{（3）和家长、同学、恋人关系出现危机，影响正常生活，怎么调整心情？}

\begin{itemize}
	\item 直面并解决问题。建议直接当面或电话沟通，线上聊天往往可能引起歧义。
	\item 吐槽。跟信得过的老师、好友吐槽。大学新增的方法是树洞吐槽。
	\item 把事情扔一边先冷静。如果不是很急于解决，先不管这件事情，正常学习生活一段时间，冷静之后再回想。
	\item 咨询。当影响很大时，千万不要藏着掖着，要敢于寻求专业的心理援助。事实上，心理咨询中心不能看作是“治病”的地方，更多时候它是一个可以放心跟人吐槽，同时有人帮你理性分析、解决问题的地方。
\end{itemize}

\textbf{（4）社恐，想社交但实在不会、不敢，怎么办？}

社交是大学的必修课。首先，不要用标签框定自己（如MBTI），要知道社交属性到大学是会变化的（例如有的同学由i转e）。其次，不同性格的人各有舒适区，具体不再赘述，请大家自行感受，并学会悦纳自我。最后，还要意识到一件事情是，社恐和社牛并没有明显的界限划分，大多数人其实是处于中间的灰色地带的。

\vspace{10pt}

认识到以上几点之后，相信“社恐”的同学就不会再过多产生社交焦虑了。如果实在觉得自己需要社交的机会，不妨先从大胆参加活动、多多认识志同道合的朋友开始吧。当然，最终能不能成为朋友很多时候是看缘分，强求不来。

\vspace{10pt}

最后提醒大家社交把握尺度，拥有自己的判断力，不要让盲目社交影响了你健康的生活和正常的自我发展。
